% Generated draft: requires mathematical and source review.
\documentclass[11pt]{article}
\usepackage[margin=25mm]{geometry}
\usepackage{fontspec}
\setmainfont{Latin Modern Roman}
\usepackage{amsmath,amssymb,mathtools,unicode-math}
\setmathfont{Latin Modern Math}
\usepackage{enumitem,needspace,xurl,hyperref}
\setlength{\parindent}{0pt}
\setlength{\parskip}{5pt}
\setlength{\emergencystretch}{3em}
\newcommand{\N}{\mathbb{N}}
\newcommand{\Z}{\mathbb{Z}}
\newcommand{\Q}{\mathbb{Q}}
\newcommand{\R}{\mathbb{R}}
\newcommand{\C}{\mathbb{C}}
\newcommand{\F}{\mathbb{F}}
\newcommand{\A}{\mathbb{A}}
\newcommand{\PP}{\mathbb{P}}
\newcommand{\E}{\mathbb{E}}
\newcommand{\eps}{\varepsilon}
\newcommand{\dd}{\,\mathrm{d}}
\begin{document}
\section*{UnsolvedMath ID: 30000044}
\textbf{Source problem code:} catalogue entry
\subsection{Definitions and mathematical statement}
For a measurable set $A\subset\R$ with $0<|A|<\infty$, equip $L^2(A)$ with Lebesgue measure and set $e_\gamma(x)=e^{2\pi\mathrm i\gamma x}$ on $A$. An unweighted exponential frame indexed by a countable discrete set $\Gamma\subset\R$ means that there are $0<a\leq b<\infty$ such that
\[
 a\|f\|_{L^2(A)}^2\leq
 \sum_{\gamma\in\Gamma}\left|\int_A f(x)e^{-2\pi\mathrm i\gamma x}\,dx\right|^2
 \leq b\|f\|_{L^2(A)}^2
 \qquad\forall f\in L^2(A).
\]
The frame is tight if $a=b$ can be chosen. The two source questions, under the finite-measure interpretation supplied by the assessment, are whether every such $A$ admits an exponential frame, and whether every such $A$ admits a tight exponential frame. Frequencies and bounds may depend on $A$; boundedness of $A$ is not assumed. Weighted exponentials, an orthogonal basis, and a merely continuous frame are different demands.

The source also phrases the question as existence of $\Gamma$ and $\eta\in H$ whose translates form a frame or a tight frame. An abstract Hilbert space alone does not define translations. One precise realization of its stated equivalence is the translation-invariant space
\[
 H_A=\{g\in L^2(\R):\widehat g=0\text{ a.e. outside }A\},
 \qquad \widehat\eta=1_A,
\]
with $T_\gamma\eta(t)=\eta(t-\gamma)$. The unitary Fourier transform takes these translates to $e^{-2\pi\mathrm i\gamma x}1_A(x)$, so their frame and tight-frame properties are exactly the corresponding exponential properties with frequency set $-\Gamma$. This explains the equivalence in a specified model, without claiming that arbitrary $H$ or arbitrary cyclic vectors give the same unweighted question. The extracted sentence leaves the original $H$, action, vector restrictions and assumptions on $A$ unstated.

The supplied literature triage checked 21 August 2026 reports existence of discrete exponential frames for every finite-measure set, including unbounded ones, as proved, but says that the stronger general tight-frame demand is not resolved. These strength levels and their dated supplied status must remain separate.
\subsection{Short English statement}
Does every positive finite-measure subset of the real line admit a discrete exponential frame, and can one always choose it tight? The supplied assessment reports the ordinary-frame assertion proved, while retaining the general tight-frame question. The translation formulation needs a specified Hilbert-space representation.
\subsection{Sources}
\begin{itemize}
\item Local UnsolvedMath registry, permanent ID 30000044, exact source code catalogue entry; \emph{Exponential Frames for Measurable Sets}. Proposed year 2004, status partially\_solved; corrected\_verified extraction preserves both frame strengths and the asserted equivalence.
\item \emph{Mini-Workshop: Wavelets and Frames (2004)}, Oberwolfach Reports, DOI 10.4171/owr/2004/10, pp. 479--544, \url{https://doi.org/10.4171/owr/2004/10}.
\item Supplied literature triage checked 2026-08-21: ``Exponential frames for unbounded sets'', \url{https://arxiv.org/abs/1410.5693}; ``Proceedings of the AMS version'', \url{https://doi.org/10.1090/proc/12868}. This dated OWR-derived curation requests independent verification of time-sensitive status and cited claims.
\end{itemize}
\end{document}
